% ============================================================
% Introduction — the four-move structure reviewers expect.
% Keep the funnel tight: every sentence moves from the broad
% problem toward YOUR contribution. 1–1.5 pages for a conference
% paper. See references/writing-guide.md before writing.
% ============================================================
\section{Introduction}
\label{sec:introduction}

% ---- Move 1: Problem context (1–2 paragraphs) ----
% Open with the broad problem and why it matters, then narrow
% quickly to the specific challenge this paper addresses. State
% the problem concretely enough that a reader can tell whether
% it affects them. Avoid empty openers ("In recent years, X has
% attracted much attention") — open with the tension instead.
LaTeX is a powerful typesetting system \cite{example2023}. For more information about conference proceedings, see \cite{conference2022}.

% ---- Move 2: Gap in existing approaches (1 paragraph) ----
% What do current solutions do, and where do they fall short?
% This paragraph justifies the paper's existence. Be specific
% about the limitation ("method X scales quadratically with
% corpus size"), not dismissive ("X is inefficient").

% ---- Move 3: This paper (1–2 paragraphs) ----
% "In this paper, we propose ..." State the key insight in one
% sentence before the mechanism. Then summarize the approach at
% a level from which a reader can reconstruct the system shape.

% ---- Contributions (itemized, verifiable) ----
% Each bullet = one claim a reviewer can verify somewhere in the
% paper. Attach the headline number to the evaluation bullet.
Our contributions are as follows:
\begin{itemize}
    \item We identify ... % the problem/insight
    \item We propose ... % the method
    \item We evaluate ... on ..., showing ... % the evidence + headline result
\end{itemize}

% ---- Organization (optional, 1 short paragraph) ----
% "The remainder of this paper is organized as follows.
% Section~\ref{sec:background} reviews background and related
% work. Section~\ref{sec:methodology} presents our approach.
% Section~\ref{sec:evaluation} reports experimental results.
% Section~\ref{sec:conclusion} concludes the paper."
